\section{Related work}
\label{sec:related}

\paragraph{Source ownership and attribution.}
Attribution evaluation across multiple candidate sources is an established measurement problem \citep{attributionbench2024}; source-aware factuality work targets cross-source conflation directly \citep{provenanceguard2026}; and scientific claim--evidence benchmarks test whether retrieved evidence supports or contradicts a claim \citep{claimbench2025}. Claim-source retrieval sharpens the same point from the retrieval side, since recovering the source a scientific statement actually came from is a different task from finding a source that would support it \citep{claimsource2026}. Together these establish claim correctness, source ownership, and semantic support as separate coordinates, which is the premise this paper builds on.

\paragraph{Provenance-sensitive authorization.}
The distinction between relevant evidence and evidence authorized to determine an action or conclusion is already established. A controlled audit holds the task and proposition fixed while varying source authority and measures whether agent actions change \citep{provsen2026}. Partial Evidence Bench studies a complementary failure mode in which access control is correct but a system overstates completeness when material evidence lies outside its authorization boundary \citep{partial2026}. AuthGraph aligns execution provenance against a separately derived authorization graph at tool and parameter-source level \citep{authgraph2026}; FAVA lowers structured task constraints into evidence-backed permission graphs checked before effectful execution \citep{fava2026}; and cryptographically verifiable authorization formalizes bindings among principal, request, context, policy, and execution evidence \citep{cva2026}. These works absorb any generic claim that provenance-aware authorization, permission graphs, or evidence-backed authorization are new here. The residual question is narrower: when the target is \emph{scientific claim promotion}, what representation and output interface make that target attainable, and does a target-bound non-compensatory relation differ from a donor-complete generic authorization relation on prospectively frozen exact contracts?

\paragraph{Auditability, citation fidelity, and influence.}
Claim-level auditability work argues that research-agent outputs require persistent semantic provenance rather than only fluent text \citep{aar2026}. ProvenAI separates citation fidelity from whether cited evidence influenced the answer path \citep{provenai2026}. We adopt that separation: citation presence cannot substitute for evidence-use provenance.

\paragraph{Research-integrity provenance and declared behavior.}
INSPECT-AI and RIPE-KG represent the provenance of research-integrity assessments over clinical-trial publications, showing a current scientific setting in which the assessment process itself is an auditable object \citep{inspectai2026}. Behavioral Integrity Verification formalizes a complementary declared-versus-actual integrity problem for agent skills \citep{behavioralintegrity2026}, and certified speculative execution shows the same instinct applied to untrusted agent actions: let the work proceed, but hold its effects behind a check that must clear before they become authoritative \citep{certifiedspeculative2026}. These are parent-domain precedents for explicit provenance and contract checking.

\paragraph{Iterative verification and evidence escalation.}
FIRE iterates retrieval and verification rather than treating a single retrieval pass as final \citep{fire2024}. DeepSciVerify motivates escalation from an abstract-level judgment to fuller evidence when initial evidence is insufficient \citep{deepsciverify2026}. Stage-wise fact-checking benchmarks make the complementary measurement argument, that an end-to-end verdict hides which stage actually failed \citep{factarena2026}. The first two mechanisms inform frozen comparator proxies here, and the third informs the decision to report per-family and per-coordinate outcomes rather than one aggregate.

\paragraph{Evaluator integrity, benchmark auditing, and contamination.}
RewardHackingAgents treats evaluator manipulation and held-out leakage as an explicit benchmark dimension \citep{rewardhacking2026}. Search-time contamination work shows that browsing systems can retrieve public benchmark answers during inference \citep{stc2026}. BenchGuard and Automated Benchmark Auditing show that benchmark specifications, environments, ground truths, and grading logic can themselves be defective \citep{benchguard2026,autobenchmarkaudit2026}. The Holistic Agent Leaderboard further demonstrates the value of standardized harnessing and trace inspection, including direct observation of agents searching for benchmark material rather than solving the intended task \citep{hal2025}. These results motivate the protected evaluator identity, holdout custody, access telemetry, and independent audit used here.

\paragraph{Abstention and scientific integrity.}
SciIntegrity-Bench constructs academic-integrity dilemmas in which honest acknowledgement of inability is the only correct outcome, while AgentAbstain evaluates paired should-act/should-abstain tasks in executable environments \citep{sciintegrity2026,agentabstain2026}. They establish abstention as an evaluation target distinct from generic task success. The object measured here is narrower: an undetermined or blocked terminal is authority-specific, produced by an unresolved or failed evidence, checker, or evaluator prerequisite rather than by a general judgement that the task should not be attempted.

